\chapter{RESULTS}
\label{ch:results}

% NOTE: This chapter framework is complete. Run ./generate.sh and test workflows,
% then fill in the TODO placeholders with actual measurements and screenshots.

This chapter presents the evaluation of the implemented IU-MiCert system through performance benchmarks, security analysis, cost analysis, and comparison with existing solutions. The chapter includes detailed happy case scenarios demonstrating complete workflows for credential issuance, student verification with selective disclosure, and credential revocation. Comprehensive analysis discusses performance achievements, usability evaluation from multiple stakeholder perspectives, security validation of the deployed system, and assessment of how the implementation meets the design goals established in Chapter 3.

\section{Performance Benchmarks}
\label{sec:performance-benchmarks}

This section presents performance measurements for the implemented system, evaluating proof sizes, timing characteristics, and efficiency comparisons with Merkle tree approaches.

\subsection{Proof Size Analysis}
\label{subsec:proof-size}

Table \ref{tab:proof-sizes} presents the measured proof sizes for IU-MiCert credentials:

\begin{table}[H]
\centering
\begin{tabular}{|p{6cm}|p{6cm}|}
\hline
\textbf{Metric} & \textbf{Value} \\
\hline
Single course proof & % TODO: Measure actual proof size (target: 600-800 bytes)
\\
\hline
Full receipt (20 courses) & % TODO: Measure full receipt size (target: 15-20 KB)
\\
\hline
On-chain storage per term & 32 bytes (root commitment only) \\
\hline
\end{tabular}
\caption{Proof Size Measurements}
\label{tab:proof-sizes}
\end{table}

% TODO: Add paragraph analyzing whether target <1KB per course was achieved
% TODO: Discuss implications of proof size for receipt portability

\subsection{Timing Benchmarks}
\label{subsec:timing-benchmarks}

Performance timing measurements evaluate the system's responsiveness and scalability:

\begin{table}[H]
\centering
\begin{tabular}{|p{7cm}|p{5cm}|}
\hline
\textbf{Operation} & \textbf{Time} \\
\hline
Tree insertion (per course) & % TODO: Measure insertion time
\\
\hline
Tree commitment (1000 courses) & % TODO: Measure commitment time
\\
\hline
Proof generation (per course) & % TODO: Measure proof generation time
\\
\hline
Proof verification (per course) & % TODO: Measure verification time
\\
\hline
Full receipt verification (20 courses) & % TODO: Measure full verification time
\\
\hline
\end{tabular}
\caption{Timing Benchmarks}
\label{tab:timing-benchmarks}
\end{table}

% TODO: Compare measured times against NFR2 target (<100ms per course)
% TODO: Discuss whether verification is fast enough for real-time applications

\subsection{Comparison with Merkle Trees}
\label{subsec:merkle-comparison}

Theoretical comparison between Merkle and Verkle tree approaches for credential systems:

\begin{table}[H]
\centering
\begin{tabular}{|p{4cm}|p{4cm}|p{4cm}|}
\hline
\textbf{Metric} & \textbf{Merkle (1M items)} & \textbf{Verkle (1M items)} \\
\hline
Proof size & 32 bytes × 20 levels = 640 bytes & Approximately 600 bytes (constant) \\
\hline
Verification complexity & O(20 hash operations) & O(1 IPA verification) \\
\hline
Tree height & 20 levels & 3 levels \\
\hline
\end{tabular}
\caption{Merkle vs Verkle Tree Comparison}
\label{tab:merkle-verkle-comparison}
\end{table}

The comparison demonstrates Verkle trees' constant proof size advantage, particularly significant as credential datasets scale to institutional sizes.

\section{Security Analysis}
\label{sec:security-analysis}

This section evaluates the security properties of the implemented system.

\subsection{Threat Model}
\label{subsec:threat-model}

Table \ref{tab:threat-model} analyzes potential threats and their mitigations:

\begin{table}[H]
\centering
\begin{tabular}{|p{5cm}|p{7cm}|}
\hline
\textbf{Threat} & \textbf{Mitigation} \\
\hline
Forge credential & IPA proof cryptographically binds course data to blockchain-anchored root \\
\hline
Modify grade after issuance & Value hash mismatch detection in verification prevents acceptance of altered data \\
\hline
Backdate credential & Blockchain timestamp is immutable; term-based trees prevent retroactive insertion \\
\hline
Replay attack with old credentials & Version tracking alerts verifiers to superseded credentials \\
\hline
Private key leakage & Only issuer can publish roots; access control in smart contract enforces authorization \\
\hline
\end{tabular}
\caption{Threat Model and Mitigations}
\label{tab:threat-model}
\end{table}

% TODO: After security testing, add results of attempted attacks (if any)

\subsection{Cryptographic Assumptions}
\label{subsec:crypto-assumptions}

The system's security relies on standard cryptographic assumptions:

\begin{table}[H]
\centering
\begin{tabular}{|p{5.5cm}|p{6.5cm}|}
\hline
\textbf{Assumption} & \textbf{Basis} \\
\hline
Discrete logarithm hardness & Bandersnatch elliptic curve used in Verkle trees \\
\hline
Hash collision resistance & SHA-256 for key and value hashing \\
\hline
Blockchain immutability & Ethereum Proof-of-Stake consensus mechanism \\
\hline
\end{tabular}
\caption{Cryptographic Security Assumptions}
\label{tab:crypto-assumptions}
\end{table}

\section{Cost Analysis}
\label{sec:cost-analysis}

This section evaluates the economic viability of blockchain-based credential systems.

\subsection{On-Chain Costs}
\label{subsec:onchain-costs}

Measured blockchain transaction costs for term root publication:

\begin{table}[H]
\centering
\begin{tabular}{|p{5.5cm}|p{3cm}|p{3.5cm}|}
\hline
\textbf{Operation} & \textbf{Gas} & \textbf{Cost (20 gwei)} \\
\hline
Publish one term root & % TODO: Get actual gas from Etherscan transaction
Approx. 50,000 & Approximately 0.001 ETH \\
\hline
Publish 7 terms (full degree) & Approximately 350,000 & Approximately 0.007 ETH \\
\hline
Read root (view function) & 0 & Free \\
\hline
\end{tabular}
\caption{On-Chain Cost Measurements}
\label{tab:onchain-costs}
\end{table}

% TODO: Update gas measurements from actual Sepolia transactions
% TODO: Add Etherscan transaction links as references

\subsection{Per-Student Cost}
\label{subsec:per-student-cost}

Cost analysis for a complete 4-year degree program:

\textbf{Assumptions:}
\begin{itemize}
    \item 7 academic terms over 4-year degree
    \item Approximately 0.001 ETH per term root publication
    \item ETH price: \$2,000 (current market rate)
\end{itemize}

\textbf{Calculation:}
\begin{equation}
\text{Total on-chain cost} = 7 \text{ terms} \times 0.001 \text{ ETH} = 0.007 \text{ ETH}
\end{equation}

\begin{equation}
\text{Cost in USD} = 0.007 \text{ ETH} \times \$2,000/\text{ETH} = \$14 \text{ total}
\end{equation}

For an institution with 1,000 students:
\begin{equation}
\text{Per-student cost} = \$14 \div 1,000 = \$0.014 \text{ per student}
\end{equation}

\textbf{Comparison with Traditional Systems:} Traditional credential verification costs \$50-\$100 per request. The IU-MiCert system reaches break-even after 0.14-0.28 verifications per student. Given that graduates typically need multiple verifications throughout careers, blockchain provides substantial cost savings.

\section{Comparison with Existing Solutions}
\label{sec:solution-comparison}

Table \ref{tab:solution-comparison} compares IU-MiCert with prominent existing blockchain credential systems:

\begin{table}[H]
\centering
\begin{tabular}{|p{3.5cm}|p{2.5cm}|p{3cm}|p{3cm}|}
\hline
\textbf{Feature} & \textbf{BlockCerts} & \textbf{IU-TransCert} & \textbf{IU-MiCert} \\
\hline
Granularity & Whole cert & Whole transcript & Per course \\
\hline
Proof size & Approx. 2KB & Approx. 1KB & Approx. 600B \\
\hline
Selective disclosure & No & No & Yes \\
\hline
Temporal proof & Limited & No & Yes (timestamps) \\
\hline
On-chain storage & Full hash & Full hash & Root only (32B) \\
\hline
Revocation & Separate list & None & Versioning \\
\hline
\end{tabular}
\caption{Feature Comparison with Existing Solutions}
\label{tab:solution-comparison}
\end{table}

IU-MiCert advances the state of blockchain credentials through micro-credential granularity, compact constant-size proofs, and elegant version-based revocation.

\section{Usability Evaluation}
\label{sec:usability}

This section evaluates the system's usability from the perspective of all stakeholder groups.

\subsection{User Flows Tested}
\label{subsec:user-flows}

% TODO: After running test workflows, update completion rates

\begin{table}[H]
\centering
\begin{tabular}{|p{3.5cm}|p{5cm}|p{3.5cm}|}
\hline
\textbf{Actor} & \textbf{Task} & \textbf{Completion Rate} \\
\hline
Issuer & Publish term root & % TODO: Test and document
\\
\hline
Student & View/filter receipt & % TODO: Test and document
\\
\hline
Verifier & Verify receipt & % TODO: Test and document
\\
\hline
Issuer & Revoke credential & % TODO: Test and document
\\
\hline
\end{tabular}
\caption{User Flow Testing Results}
\label{tab:user-flows}
\end{table}

\subsection{Interface Screenshots}
\label{subsec:interface-screenshots}

This section presents the implemented user interfaces for all system components.

\subsubsection{Issuer Dashboard Interface}

\begin{figure}[H]
\centering
\fbox{\parbox{0.9\textwidth}{
\centering
\vspace{1cm}
\textbf{[Screenshot 6.1: Issuer Dashboard - Main Page]}\\[0.5cm]
\textit{Description: Main dashboard showing term management interface with options to publish term roots, generate demo data, and manage revocations. Displays list of terms with publication status (unpublished/published), term root hashes, and action buttons.}\\[0.5cm]
\textit{(Placeholder - screenshot to be captured before final submission)}
\vspace{1cm}
}}
\caption{Issuer Dashboard Main Interface}
\label{fig:screenshot61}
\end{figure}

% TODO: Replace placeholder with actual screenshot from https://iumicert-issuer.vercel.app

\begin{figure}[H]
\centering
\fbox{\parbox{0.9\textwidth}{
\centering
\vspace{1cm}
\textbf{[Screenshot 6.2: Issuer Dashboard - Publish Term Root]}\\[0.5cm]
\textit{Description: Term root publication interface showing the term identifier, computed Verkle root (32-byte hash), blockchain transaction status, and confirmation dialog. Displays gas estimation and Sepolia testnet connection indicator.}\\[0.5cm]
\textit{(Placeholder - screenshot to be captured before final submission)}
\vspace{1cm}
}}
\caption{Term Root Publication Interface}
\label{fig:screenshot62}
\end{figure}

\begin{figure}[H]
\centering
\fbox{\parbox{0.9\textwidth}{
\centering
\vspace{1cm}
\textbf{[Screenshot 6.3: Issuer Dashboard - Revocation Management]}\\[0.5cm]
\textit{Description: Revocation management page displaying pending revocation requests table with columns for student ID, course ID, term, reason, and action buttons (approve/reject). Shows version history for affected terms.}\\[0.5cm]
\textit{(Placeholder - screenshot to be captured before final submission)}
\vspace{1cm}
}}
\caption{Revocation Management Interface}
\label{fig:screenshot63}
\end{figure}

\subsubsection{Student/Verifier Portal Interface}

\begin{figure}[H]
\centering
\fbox{\parbox{0.9\textwidth}{
\centering
\vspace{1cm}
\textbf{[Screenshot 6.4: Student Portal - Receipt Upload]}\\[0.5cm]
\textit{Description: Clean interface with drag-and-drop zone for JSON receipt upload. Shows file validation status and preview of receipt contents before verification.}\\[0.5cm]
\textit{(Placeholder - screenshot to be captured before final submission)}
\vspace{1cm}
}}
\caption{Receipt Upload Interface}
\label{fig:screenshot64}
\end{figure}

\begin{figure}[H]
\centering
\fbox{\parbox{0.9\textwidth}{
\centering
\vspace{1cm}
\textbf{[Screenshot 6.5: Verification Results - Success]}\\[0.5cm]
\textit{Description: Successful verification result page showing green checkmark, verified courses list with details (course ID, name, grade, credits), blockchain confirmation status, and term root validation details.}\\[0.5cm]
\textit{(Placeholder - screenshot to be captured before final submission)}
\vspace{1cm}
}}
\caption{Successful Verification Results}
\label{fig:screenshot65}
\end{figure}

\begin{figure}[H]
\centering
\fbox{\parbox{0.9\textwidth}{
\centering
\vspace{1cm}
\textbf{[Screenshot 6.6: Verification Results - Detailed View]}\\[0.5cm]
\textit{Description: Expanded view showing per-course verification details including proof size, verification time, blockchain query results, and term version information.}\\[0.5cm]
\textit{(Placeholder - screenshot to be captured before final submission)}
\vspace{1cm}
}}
\caption{Detailed Verification View}
\label{fig:screenshot66}
\end{figure}

\section{Happy Case Scenarios}
\label{sec:happy-cases}

% NOTE: Run ./generate.sh to create receipts, then execute these workflows
% and capture actual measurements to fill in the placeholders below

This section demonstrates the system through complete user flows showing successful operations from start to finish. Each scenario presents a realistic use case with detailed steps, measurements, and outcomes.

\subsection{Scenario 1: Complete Credential Issuance}
\label{subsec:scenario-issuance}

\textbf{Context:} University completes an academic term and needs to issue verifiable credentials for all students.

\textbf{Steps:}

\textbf{1. Data Preparation}

% TODO: Document actual data preparation process
\begin{itemize}
    \item Issuer exports final grades from LMS as JSON file
    \item File contains: % TODO: Document actual number of students and courses
    \item Total course completions: % TODO: Calculate total
\end{itemize}

\textbf{2. Tree Construction} (See Figure \ref{fig:screenshot67})

\begin{figure}[H]
\centering
\fbox{\parbox{0.9\textwidth}{
\centering
\vspace{1cm}
\textbf{[Screenshot 6.7: Tree Construction Progress]}\\[0.5cm]
\textit{Description: Dashboard showing tree construction in progress with progress bar, number of courses inserted, and estimated completion time.}\\[0.5cm]
\textit{(Placeholder - screenshot to be captured during actual workflow execution)}
\vspace{1cm}
}}
\caption{Tree Construction Progress}
\label{fig:screenshot67}
\end{figure}

% TODO: Run tree construction and measure:
\begin{itemize}
    \item Tree construction time: % TODO: Measure
    \item Computed root: % TODO: Record actual root hash
\end{itemize}

\textbf{3. Root Publication} (See Figure \ref{fig:screenshot68})

\begin{figure}[H]
\centering
\fbox{\parbox{0.9\textwidth}{
\centering
\vspace{1cm}
\textbf{[Screenshot 6.8: Blockchain Publication]}\\[0.5cm]
\textit{Description: Transaction submission interface showing Metamask confirmation dialog, gas estimate, and transaction status. After confirmation, shows Etherscan link to transaction.}\\[0.5cm]
\textit{(Placeholder - screenshot to be captured during actual blockchain publication)}
\vspace{1cm}
}}
\caption{Blockchain Root Publication}
\label{fig:screenshot68}
\end{figure}

% TODO: Execute publish-roots command and record:
\begin{itemize}
    \item Gas used: % TODO: Get from Etherscan
    \item Transaction hash: % TODO: Record
    \item Confirmation time: % TODO: Measure
\end{itemize}

\textbf{4. Receipt Generation} (See Figure \ref{fig:screenshot69})

\begin{figure}[H]
\centering
\fbox{\parbox{0.9\textwidth}{
\centering
\vspace{1cm}
\textbf{[Screenshot 6.9: Receipt Generation Progress]}\\[0.5cm]
\textit{Description: Dashboard showing receipt generation for all students with progress indicator, number of receipts generated, and estimated time remaining.}\\[0.5cm]
\textit{(Placeholder - screenshot to be captured during receipt generation)}
\vspace{1cm}
}}
\caption{Receipt Generation Progress}
\label{fig:screenshot69}
\end{figure}

% TODO: Run generate-all-receipts and measure:
\begin{itemize}
    \item Receipts generated: % TODO: Count
    \item Average size: % TODO: Calculate
    \item Generation time: % TODO: Measure
\end{itemize}

\textbf{Result:} % TODO: Summarize credentials verified, blockchain commitment, total cost

\subsection{Scenario 2: Student Credential Verification}
\label{subsec:scenario-verification}

\textbf{Context:} Student applies for internship requiring proof of specific programming courses.

\textbf{Steps:}

\textbf{1. Receipt Access} (See Figure \ref{fig:screenshot610})

\begin{figure}[H]
\centering
\fbox{\parbox{0.9\textwidth}{
\centering
\vspace{1cm}
\textbf{[Screenshot 6.10: Student Receipt Download]}\\[0.5cm]
\textit{(Placeholder - screenshot from student portal)}
\vspace{1cm}
}}
\caption{Student Receipt Download}
\label{fig:screenshot610}
\end{figure}

% TODO: Document: terms, courses, file size

\textbf{2. Selective Disclosure} (See Figure \ref{fig:screenshot611})

\begin{figure}[H]
\centering
\fbox{\parbox{0.9\textwidth}{
\centering
\vspace{1cm}
\textbf{[Screenshot 6.11: Receipt Filtering]}\\[0.5cm]
\textit{(Placeholder)}
\vspace{1cm}
}}
\caption{JSON Filtering}
\label{fig:screenshot611}
\end{figure}

% TODO: Document filtering

\textbf{3-5. Verification Process} (See Figures \ref{fig:screenshot612}-\ref{fig:screenshot614})

\begin{figure}[H]
\centering
\fbox{\parbox{0.9\textwidth}{\centering\vspace{1cm}\textbf{[Screenshot 6.12-6.14: Verification Flow]}\vspace{1cm}}}
\caption{Upload, Verify, Result}
\label{fig:screenshot612}
\end{figure}

% TODO: Run verification, measure times

\subsection{Scenario 3: Credential Revocation}
\label{subsec:scenario-revocation}

% TODO: Test revocation workflow

\textbf{Context:} Academic misconduct requires credential revocation.

% TODO: Document steps 1-5 with actual measurements and screenshots 6.15-6.18

\section{Analysis and Discussion}
\label{sec:results-analysis}

% TODO: After completing measurements, write analysis discussing:
% - Performance achievements vs targets
% - Usability from all perspectives
% - Security validation
% - Design goal attainment

\subsection{Performance Analysis}
\label{subsec:perf-analysis}

% TODO: Write analysis of proof sizes, verification speeds, cost efficiency

\subsection{Usability Analysis}
\label{subsec:usability-analysis}

% TODO: Analyze issuer, student, and verifier experiences

\subsection{Security Analysis}
\label{subsec:security-analysis-deployed}

% TODO: Analyze blockchain immutability, proof integrity, revocation correctness

\subsection{Design Goals Validation}
\label{subsec:design-validation}

% TODO: Validate all SR1-SR8 requirements met
